\documentclass[10pt,letterpaper]{article}
\usepackage[T1]{fontenc}
\usepackage[utf8]{inputenc}
\usepackage{lmodern}
\usepackage[margin=0.76in]{geometry}
\usepackage[table]{xcolor}
\usepackage{array,tabularx,longtable,booktabs}
\usepackage{enumitem}
\usepackage{amssymb,amsmath}
\usepackage{hyperref}
\usepackage{fancyhdr,lastpage}
\usepackage{microtype}
\usepackage{parskip}
\definecolor{cdblue}{HTML}{203957}
\definecolor{cdgray}{HTML}{586677}
\hypersetup{colorlinks=true,linkcolor=cdblue,urlcolor=cdblue,pdftitle={Chapter 02: Configuration Management}}
\pagestyle{fancy}
\fancyhf{}
\lhead{\small\color{cdgray}Continuous Delivery Engineering}
\rhead{\small\color{cdgray}Study guide}
\cfoot{\thepage\ / \pageref{LastPage}}
\setlength{\headheight}{14pt}
\setlength{\emergencystretch}{2em}
\setlist[itemize]{leftmargin=*,itemsep=3pt,topsep=4pt}
\setlist[enumerate]{leftmargin=*,itemsep=3pt,topsep=4pt}
\begin{document}
\begin{center}
{\LARGE\bfseries\color{cdblue}Chapter 02: Configuration Management}\\[7pt]
{\large I: Foundations}\\[6pt]
{\small Reference edition: 2026-10-30 $\bullet$ Original educational and operational material}
\end{center}
\noindent\color{cdblue}\rule{\linewidth}{1.2pt}\color{black}
\vspace{0.4em}
\section{Chapter scope and provenance}
Primary reading: Humble and Farley, Continuous Delivery (2010), Chapter 2, "Configuration Management"; Part I: Foundations. This guide is an original synthesis and not a reproduction of the chapter. The book provides the foundational delivery principles; current AWS, Terraform and GitHub Actions references in the applied sections are supplementary.


\section{Learning outcomes}
\begin{enumerate}
\item Define the configuration items required to recreate a release.
\item Distinguish secrets, environment configuration, code and built artifacts.
\item Identify drift and uncontrolled production mutation.
\end{enumerate}

\section{Core conceptual model}
Focus areas: version control; software configuration; reproducible environments; external libraries.

\begin{itemize}
\item Input: a version-controlled change, explicit requirements, and a reproducible execution context.
\item Mechanism: controlled feedback, automation, evidence capture, and a stop condition where risk exceeds tolerance.
\item Output: a verifiable release candidate or a diagnosed failure with an accountable owner.
\end{itemize}

\section{Reading and analysis protocol}
\begin{enumerate}
\item Read the chapter and record five propositions, distinguishing an author claim from an example or assumption.
\item Identify two failure modes, one recovery technique, and one measurable acceptance condition.
\item Trace each technique to the stages and controls in the deployment-pipeline-architecture reference.
\item State how the proposed practice changes for a regulated system with segregated duties.
\end{enumerate}

\section{Detailed technical study notes}
\subsection{Configuration item inventory}
Track source, dependency versions, build scripts, provisioning, deployment commands, test fixtures and environment policy as distinct controlled assets.

\subsection{Change authority}
A version-controlled change should be reviewable, attributable and reversible where practical. Emergency changes require subsequent reconciliation rather than permanent drift.

\subsection{Environment parameters}
Separate immutable program logic from environment-specific inputs. Audit nonsecret parameters and protect secrets in a dedicated, access-controlled store.

\subsection{Transitive dependencies}
Pin resolved versions and track origin and integrity. A library update can change runtime behavior even when application source is unchanged.

\subsection{Drift and state}
Record expected desired state, detect deviation, and decide if changes are reconciled automatically or escalated for review.

\subsection{Configuration recovery}
Practice rebuilding an environment from controlled inputs to prove the inventory is sufficiently complete.

\section{Domain-specific challenge questions}
\begin{enumerate}
\item For configuration item inventory, identify the precondition, expected evidence, adversarial failure case and measurable corrective action.
\item For change authority, identify the precondition, expected evidence, adversarial failure case and measurable corrective action.
\item For environment parameters, identify the precondition, expected evidence, adversarial failure case and measurable corrective action.
\item For transitive dependencies, identify the precondition, expected evidence, adversarial failure case and measurable corrective action.
\end{enumerate}
\section{Guided exercise}
Build a configuration manifest and a dependency lockfile policy for dev, staging and prod.

\subsection{Required evidence}
\begin{itemize}
\item A version-controlled design or executable artifact.
\item A failure demonstration, negative test or counterexample.
\item An improvement statement with baseline, target, owner and verification method.
\end{itemize}


\section{Engineering scenario}
An engineer hot-fixes a production environment outside version control; explain recovery and preventive controls.

\begin{itemize}
\item Describe the failure propagation path and the earliest feasible detection point.
\item Recommend a minimal remediation and its rollback or containment plan.
\item Define a measure that would demonstrate risk reduction.
\end{itemize}

\section{Knowledge checks}
\begin{enumerate}
\item Which condition is the gate, and which action is the response when it fails?
\item What evidence distinguishes successful automation from an assumed successful release?
\item Where might environments, data, dependencies or people invalidate the technique?
\item What feedback could be moved earlier without losing necessary confidence?
\item Which assumptions are book principles, and which are modern supplementary implementation choices?
\end{enumerate}

\section{Completion rubric}
Score each dimension from 0 (missing) to 4 (demonstrated and independently verifiable): conceptual accuracy; reproducibility; negative testing; security and safety; documentation and traceability. Pass threshold: 16/20 with no zero on reproducibility or safety.


\section{Prerequisites and onward route}
Git, basic tests, and a build runner. Next: Chapters 3 onward.



\end{document}
